\paragraph{Three charged minima of a finite-range native source.}
The full native Higgs tensor admits a declared nearest-neighbour regulator
with charged-branch dispersion
\[
 E_k(p)^2=\frac{\sum_j\sin^2(ap_j)}{a^2}
 +\left[m_k+\frac{\sum_j(1-\cos(ap_j))}{a}\right]^2.
\]
For $x_j=ap_j$, differentiation gives
$\partial_{x_j}E_k^2=2\sin x_j
[am_k+1+\sum_{i\ne j}(1-\cos x_i)]/a^2$.
Since $m_k>0$, each charged branch has exactly one local and global
minimum on the Brillouin torus, at zero momentum.
The native colour and charge operators identify three colourless charged
Dirac branches, with electron-anchored energies
$0.00051099895069$, $0.1056594144959$ and $1.7769849708134\,\mathrm{GeV}$.
The complete side-four execution retains all $11\,520$ doubled rows,
all colourful and neutral modes, and complete creation/arrival records.
The former charged corner minima are retained eigenstates with higher
energy and negative curvature in at least one direction.
The neutral sine valleys and the degenerate neutral reference are retained.
The regulator coefficient is declared, and the analyzer phase $2/9$
has retrospective ancestry. No interacting particle identification or
chiral gauge continuum is inferred. On momenta $\lVert p\rVert\leq P$,
the operator error from the continuum Dirac branch is at most
$aP^2/2+a^2P^3/6$; complete same-input record variation is at most
$\min\{1,|t|(aP^2/2+a^2P^3/6)\}$.
